\section{Introduzione}
\label{sec:introduzione}

L'obiettivo del progetto è lo sviluppo di un'applicazione distribuita
nel compute continuum per il monitoraggio di misure ambientali, distribuendo
l'elaborazione tra livelli Edge e Cloud. Il caso di studio utilizza dati
provenienti da sensori BME280 del dataset pubblico Sensor.Community, relativi
a temperatura, umidità e pressione e raccolti nel mese di gennaio 2025.

Il dataset storico viene utilizzato per simulare un sistema di monitoraggio
continuo. Le osservazioni vengono riprodotte preservandone le distanze
temporali in event time mediante un replay accelerato e inviate ai rispettivi
Edge Site. Ciascun Edge esegue una prima fase di validazione e aggregazione
locale, mentre il livello Cloud combina successivamente i risultati provenienti
dalle diverse zone per ottenere una vista complessiva del sistema. Questa
suddivisione permette di distribuire la computazione lungo il continuum e di
ridurre il numero di record trasferiti verso il Cloud, inoltrando aggregati
locali anziché ogni singola osservazione.

Poiché il progetto è svolto individualmente, tra i requisiti alternativi
previsti dalla traccia è stata scelta la scalabilità orizzontale. La
valutazione sperimentale studia quindi il comportamento del sistema al variare
del numero di Cloud Worker, mantenendo invariato il workload, e analizza
inoltre l'impatto delle condizioni di rete sulla pipeline distribuita.

I componenti runtime dell'applicazione sono sviluppati in Go, containerizzati
tramite Docker e distribuiti su AWS mediante un processo di deployment
automatizzato. Il preprocessing offline del dataset è invece realizzato in
Python.